Manual tracking of unauthorized graffiti becomes infeasible as image banks
grow. This work designs, implements and experimentally evaluates a modular
system to cluster an unlabelled dataset of graffiti photographs by visual
similarity, aiming to characterize both clustering quality and the
computational cost of its components as a function of bank size. The
architecture chains swappable stages (segmentation, \textit{embedding}
extraction, storage, dimensionality reduction and clustering), and a
benchmarking framework sweeps combinations of them, measuring per-stage time,
memory and quality metrics for each run. The resulting base configuration (a
\textit{fine-tuned} DINOv2 style head + UMAP-10 + HDBSCAN + full image) is
validated both on the dataset of $6416$ graffiti images and on disjoint
manually labelled subsets. The style head separates four styles with an
adjusted Rand index of 0.722 (K-Means, oracle $k{=}4$), whereas author
clustering stays near chance, exposing the insufficient per-author annotations
in the test set. The scaling study shows that SQLite similarity search and
ingestion are the dominant bottlenecks at this scale, while other factors such
as the complexity of the reduction or clustering algorithms would only start to
dominate at much larger scales. We conclude that intrinsic
clustering-evaluation metrics can be misleading and that the choice of feature
extractor must be anchored in supervised evidence.
